\documentclass[12pt,a4paper]{article}

\usepackage{fontspec}
\setmainfont{Latin Modern Roman}
\newfontfamily{\hindifont}{Nirmala UI}[Script=Devanagari]
\newcommand{\hindi}[1]{{\hindifont #1}}

\usepackage[margin=2.5cm]{geometry}
\usepackage{xcolor}
\usepackage{enumitem}
\usepackage{hyperref}
\usepackage{tcolorbox}
\usepackage{booktabs}
\usepackage{tabularx}
\usepackage{fancyhdr}
\usepackage{titlesec}
\usepackage{setspace}

\tcbuselibrary{skins, breakable}

% Colours
\definecolor{govblue}{HTML}{003366}
\definecolor{govgold}{HTML}{C19A3E}
\definecolor{accent1}{HTML}{1B4F72}
\definecolor{accent2}{HTML}{2E86C1}
\definecolor{accent3}{HTML}{27AE60}

% tcolorbox styles
\newtcolorbox{frameworkbox}[1][]{
  colback=govblue!5, colframe=govblue,
  fonttitle=\bfseries, title={#1},
  boxrule=0.5pt, arc=2mm, breakable
}

% Header/Footer
\pagestyle{fancy}
\fancyhf{}
\fancyhead[L]{\small\textcolor{govblue}{AIGGPA Internship -- Concept Proposal}}
\fancyhead[R]{\small\textcolor{govblue}{April 2026}}
\fancyfoot[C]{\thepage}
\renewcommand{\headrulewidth}{0.4pt}
\setlength{\headheight}{14pt}

% Section Formatting
\titleformat{\section}
  {\Large\bfseries\color{govblue}}{\thesection}{1em}{}[\titlerule]
\titleformat{\subsection}
  {\large\bfseries\color{accent1}}{\thesubsection}{1em}{}

\setstretch{1.25}
\hypersetup{colorlinks=true, linkcolor=govblue, urlcolor=accent2}

\begin{document}

\begin{center}
\vspace*{0.5cm}
{\Huge\bfseries\textcolor{govblue}{Research Concept Proposal}}

\vspace{0.3cm}
{\Large\textcolor{accent1}{Field Assessment of HR-IT Integration \& Technology Adoption}}

{\large\textcolor{accent1}{\hindi{मानव संसाधन-सूचना प्रौद्योगिकी एकीकरण का क्षेत्रीय मूल्यांकन}}}

\vspace{0.5cm}
\rule{0.6\textwidth}{1pt}

\vspace{0.3cm}
{\large Atal Bihari Vajpayee Institute of Good Governance \& Policy Analysis (AIGGPA)}
\end{center}

\vspace{1cm}

\section{1. Background \& Purpose (\hindi{पृष्ठभूमि एवं उद्देश्य})}

The Government of Madhya Pradesh has made significant strides in digital governance through initiatives anchored in the principles of Mission Karmayogi (\hindi{मिशन कर्मयोगी}). However, secondary research at AIGGPA reveals an explicit \textbf{data gap}: there is no existing internal dataset or research paper tracking the actual, ground-level integration of Human Resource and Information Technology (HR-IT) tools among state and district officials.

The purpose of this research is to conduct a primary, exploratory study to measure unbiased technology adoption levels, identify capacity-building bottlenecks, and categorize officials into a Three-Tier Awareness Framework (Active Users, Aware Non-Users, Unaware).

\section{2. The Government's Role \& Operational Matrix (\hindi{सरकार की भूमिका})}

The shift from a "Rule-Based" to a "Role-Based" bureaucracy requires officials to perform their duties efficiently while acting as agile facilitators of public welfare. This research evaluates technology adoption across two distinct operational dimensions:

\begin{frameworkbox}[The Technology Dimensions]
\begin{enumerate}
    \item \textbf{Technology used to serve the public (G2C):} Portals, applications, and dashboards used for scheme delivery, grievance redressal, and citizen interaction.
    \item \textbf{Technology used for personal administrative use (G2G):} e-Office, e-HRMS, SPARROW, and internal asset/file management systems that enhance the official's personal efficiency.
\end{enumerate}
\end{frameworkbox}

\section{3. Technology Audit \& Scope (\hindi{प्रौद्योगिकी ऑडिट एवं दायरा})}

To ensure statistically significant and actionable findings, this audit will target three specific, high-impact sectors spanning both State and District level offices:
\begin{itemize}
    \item \textbf{Education Department (\hindi{शिक्षा विभाग})}
    \item \textbf{Health Department (\hindi{स्वास्थ्य विभाग})}
    \item \textbf{Forest Department (\hindi{वन विभाग})}
\end{itemize}

The audit will seek answers to critical operational questions: Do employees use e-Office? Is archiving and digitisation occurring? Do employees require additional tools (e.g., Google Workspace/Microsoft equivalents)? What legacy/parallel tech are they using instead?

\section{4. Research Methodology (\hindi{अनुसंधान पद्धति})}

\begin{itemize}
    \item \textbf{Approach:} Primary Exploratory Research.
    \item \textbf{Instrument:} A structured, bi-lingual (English/Hindi) survey designed to elicit unbiased opinions, avoiding leading questions.
    \item \textbf{Execution:} Direct field interactions at district/state offices. To combat "survey fatigue" and ensure honest responses, the survey will be presented as an efficiency audit rather than an individual performance metric.
    \item \textbf{Respondents:} Stratified sampling across Class I, II, III, and IV officials to determine where the technical skill gap is widest.
\end{itemize}

\section{5. Deliverables \& Capacity Building Solutions}

The final output will bridge the identified data gap at AIGGPA, resulting in:
\begin{enumerate}
    \item A quantitative \textbf{Adoption Index Matrix} plotting usage vs. awareness.
    \item A \textbf{Capacity Building Action Plan} tailored for the three target departments.
    \item A policy advisory on streamlining internal file management and asset tracking.
\end{enumerate}

\end{document}
